\section{Detectors}
\label{sec:detectors}

One subsection per hazard. Each names the people whose work it uses, states
what we changed, and reports only figures we measured ourselves; figures
published by the original authors are listed separately and attributed.

Five hazards are integrated. Four are disaster hazards and carry the system's
purpose; the fifth, road incidents, is retained explicitly as a generality
experiment for the validator rather than as a headline capability — a car
accident is an emergency, not a disaster. Its value is that it has completely
different physics and timescales from wildfire, so a validator that works on
both is demonstrably not tuned to one phenomenon.

\subsection{Fire and smoke}
\label{sec:fire}

\begin{description}[leftmargin=7.4em,labelwidth=7.0em,style=sameline,font=\normalfont\itshape]
  \item[Model] \texttt{yolo11s\_sensitive-detector}, Apache-2.0.
  \item[Authors] The PyroNear association, a French open-source non-profit, in
    collaboration with the French Fire and Rescue Services
    (SDIS)~\citep{pyronear2025}.
  \item[Deployment] 15 lookout towers, 51 cameras across France, Spain and
    Chile, running on Raspberry Pi hardware.
  \item[Envelope] Fixed-mount wide-field outdoor cameras viewing a landscape
    horizon, in daylight. Not validated for indoor fire, close-range flame or
    night operation.
\end{description}

It is a deliberately \emph{sensitive} detector: the authors recommend a low
confidence floor ($0.20$) and a very low NMS $\iou$ ($0.01$), trading
precision for recall on the assumption that a downstream filter exists. That
design choice is precisely why it is the right baseline for this project — the
filter it assumes is what we built.

We run the published ONNX export unmodified. The model is downloaded from its
original distributor at setup time and is never redistributed here. We
relabelled its single class from \texttt{item} to \texttt{smoke}, and wrote
our own letterbox preprocessing and YOLO decoding with NumPy non-maximum
suppression, so the runtime has no dependency on the AGPL-licensed package
(\cref{sec:licensing}).

\begin{table}[htbp]
\centering\small
\caption{Fire detector, single-frame and with no temporal validation — this is
the baseline the validator must improve on, not a headline result.
\texttt{pyro-sdis}~\citep{pyrosdis} validation split, \FireImages{} rows
(\FirePositives{} positive, \FireNegatives{} negative, 21 cameras), confidence
$0.20$, $\tau = 0.3$. \measured{} — \texttt{eval/run\_image\_eval.py}.}
\label{tab:fire}
\begin{tabular}{l S[table-format=1.3] l}
\toprule
Metric & {Value} & Note \\
\midrule
Box precision              & \FireBoxPrecision & \\
Box recall                 & \FireBoxRecall    & \\
Box $F_2$                  & \FireBoxFTwo      & recall-weighted, \cref{eq:fbeta} \\
Image-level recall         & \FireImageRecall  & \\
Image-level FAR            & \FireImageFAR     & 95\,\% CI \FireImageFARCI \\
Median latency, ms & \FireLatency & M4 Pro, CPU provider \\
\bottomrule
\end{tabular}
\end{table}

The authors publish no precision, recall or mAP for this checkpoint, so
\cref{tab:fire} has no \reported{} column: there is nothing to compare
against, and inventing one would be worse than the gap.

\begin{recorded}{A caveat that earned its keep}
The most quoted figure in this project — that the published fire detector
alarms on frames containing no smoke — was first measured on
\FireNarrowNegatives{} negative frames and carried the standing note
``indicative, not precise; widen before publishing''. Widening it to
\FireNegatives{} negatives from 21 cameras, a twenty-five-fold increase, moved
the false-alarm rate from \FireNarrowFAR{} to \FireImageFAR{} and box
precision from \FireNarrowBoxPrecision{} to \FireBoxPrecision. Recall barely
moved (\FireNarrowBoxRecall{} to \FireBoxRecall), which is the reassuring
half: the detector finds what it finds, and it was the small sample that had
flattered it on false alarms. The motivating claim is untouched — a detector
alarming on 71\,\% of clear frames is unusable without a filter — but the
specific number was overstated by sixteen percentage points and every
quotation of it was corrected. \Cref{tab:intervals} shows why the original was
never quotable: its interval was \FireNarrowFARCI.
\end{recorded}

\begin{figure}[htbp]
\centering
\includegraphics[width=0.85\textwidth]{fire_negative_widening.pdf}
\caption{The same detector at the same threshold, measured against
\FireNarrowNegatives{} negatives and then against \FireNegatives{}. Recall
barely moved; precision and the false-alarm rate both fell. This is what a
small sample flattering a detector looks like.}
\label{fig:widening}
\end{figure}

\subsection{Flood and water level}
\label{sec:flood}

The keystone hazard, and the one directly connected to the DANA. Two
capabilities: segmenting the water region, and estimating water level against
a fixed per-camera reference line (\cref{sec:flood-trend}).

\begin{description}[leftmargin=7.4em,labelwidth=7.0em,style=sameline,font=\normalfont\itshape]
  \item[Architecture] DeepLabV3~\citep{chen2017rethinking} with a
    MobileNetV3-Large backbone~\citep{howard2019searching}, from torchvision
    (BSD-3-Clause).
  \item[Training data] ATLANTIS~\citep{erfani2022atlantis} — 5{,}195
    pixel-annotated waterbody images, 56 classes.
  \item[Trained by] Us. This is the only ground-level detector in VIGÍA that
    we trained ourselves.
  \item[Envelope] Water region segmentation in outdoor scenes, ground and
    oblique viewpoints. The per-camera reference line is optional: without
    calibration the uncalibrated area signal still yields a trend, so a camera
    is useful the moment it is connected.
\end{description}

The architecture was chosen for licensing as much as performance: torchvision
is BSD-3-Clause, making this the only detector in the system entirely free of
AGPL. ATLANTIS's 56 classes were collapsed to binary water by
\cref{eq:merge} — VIGÍA does not need to know whether it is looking at a canal
or a river, only where the water is and whether it is rising. Seventeen
liquid-water labels count as water; water-associated \emph{structures} (dam,
levee, pier, culvert, breakwater) are excluded because they are not water, as
are glaciers and snow, which are not liquid, and mangrove, which is vegetation
standing in water whose extent does not track a level. The class mapping was
derived empirically by cross-referencing the dominant mask value in each label
folder against the published label list, not assumed.

\begin{table}[htbp]
\centering\small
\caption{Flood segmenter on the ATLANTIS test split, \FloodTestImages{}
images, held out and never seen during training or validation. No third-party
baseline is quoted because the dataset authors' own network reports over all
56 classes rather than binary water, so its figures are not comparable.
\measured{} — \texttt{eval/run\_flood\_eval.py}.}
\label{tab:flood}
\begin{tabular}{l S[table-format=1.4] l}
\toprule
Metric & {Value} & Note \\
\midrule
Water $\iou$    & \FloodWaterIoU      & pooled, \cref{eq:iou-pooled} \\
Precision       & \FloodPrecision     & \\
Recall          & \FloodRecall        & \\
$F_1$           & \FloodFOne          & \\
Pixel accuracy  & \FloodPixelAccuracy & \\
Median latency, ms & \FloodLatency & M4 Pro, CPU provider \\
\bottomrule
\end{tabular}
\end{table}

\begin{recorded}{A deployment hazard in the export}
PyTorch's ONNX exporter splits weights into a separate \texttt{.onnx.data}
sidecar. Copying only the \texttt{.onnx} yields a graph that loads
successfully and then produces garbage. Our export folds any sidecar back into
the single file and verifies argmax parity against the source PyTorch model,
failing below 99.9\,\% agreement.
\end{recorded}

\subsubsection{The aerial variant}

A UAV camera declaring a \textsc{nadir} viewpoint was previously refused by
\cref{def:admissible}, because the ATLANTIS model's envelope excludes it. A
second segmenter was therefore trained on FloodNet Track
1~\citep{rahnemoonfar2021floodnet} (CDLA-Permissive-1.0) and the pipeline
routes by viewpoint.

Its evidence class is weaker than everything else in this document and is
labelled as such: FloodNet withheld its validation and test masks for the
competition, so no independent labelled split exists in the release. The
figures below come from a stratified split carved out of the
\AerialLabelled{} labelled training pairs and \emph{seen by model selection} —
training-time validation, not held-out evaluation.

\begin{table}[htbp]
\centering\small
\caption{Aerial flood segmenter. Two $\iou$ figures are reported because
\cref{sec:pooling} requires it: only \AerialFlooded{} of \AerialLabelled{}
labelled images are flooded, so the pooled figure is carried by dry scenes.
Model selection used the stratified figure. \textsc{validation-measured}, not
held out — \texttt{scripts/train\_flood\_aerial.py}.}
\label{tab:flood-aerial}
\begin{tabular}{l S[table-format=1.4] l}
\toprule
Metric & {Value} & Note \\
\midrule
Water $\iou$, flooded subset & \AerialFloodedIoU   & \cref{eq:iou-strat}; the honest figure \\
Water $\iou$, pooled         & \AerialOverallIoU   & \cref{eq:iou-pooled}; flattered by dry scenes \\
Pixel accuracy               & \AerialPixelAccuracy & \\
Best epoch                   & \AerialBestEpoch    & selected on the flooded subset \\
\bottomrule
\end{tabular}
\end{table}

\begin{figure}[htbp]
\centering
\includegraphics[width=0.85\textwidth]{aerial_flood_training.pdf}
\caption{Training the aerial flood segmenter. The two curves are reported
separately because only \AerialFlooded{} of \AerialLabelled{} labelled images
are flooded; a single pooled $\iou$ would be carried by dry scenes and would
flatter the model on the one case it exists to handle.}
\label{fig:aerial}
\end{figure}

\subsection{People in water}
\label{sec:piw}

\begin{description}[leftmargin=7.4em,labelwidth=7.0em,style=sameline,font=\normalfont\itshape]
  \item[Model] RF-DETR Small fine-tuned on SeaDronesSee~\citep{rfdetrseadronessee},
    Apache-2.0.
  \item[Dataset] SeaDronesSee~\citep{varga2022seadronessee}, CC0 1.0 — the
    least restrictive licence in the project, so the credit here is given by
    choice rather than obligation.
  \item[Why RF-DETR] The same author published YOLO11n and YOLO11x checkpoints
    for this dataset under AGPL-3.0. RF-DETR was chosen specifically because
    it is Apache-2.0 and keeps this detector outside the AGPL perimeter.
  \item[Envelope] Elevated and aerial views of open water,
    \SIrange{5}{260}{\metre} altitude. Not validated for pool settings, surf
    zones or extreme crowd density.
\end{description}

Only the \texttt{swimmer} class reaches the validator. A boat is context, not
an emergency. RF-DETR emits a fixed set of object queries with no duplicate
boxes, so unlike the YOLO detectors it requires no non-maximum suppression.

\begin{table}[htbp]
\centering\small
\caption{People-in-water detector, SeaDronesSee validation split,
\SwimmerImages{} images (\SwimmerPositiveImages{} with swimmers),
\SwimmerInstances{} ground-truth swimmer instances, confidence $0.30$.
\measured{} — \texttt{eval/run\_person\_in\_water\_eval.py}.}
\label{tab:piw}
\begin{tabular}{l S[table-format=1.4] l}
\toprule
Metric & {Value} & Note \\
\midrule
Swimmer box precision & \SwimmerPrecision & \\
Swimmer box recall    & \SwimmerRecall    & 95\,\% CI \SwimmerBoxRecallCI \\
Swimmer box $F_1$     & \SwimmerFOne      & \\
Image-level FAR       & \SwimmerImageFAR  & 95\,\% CI \SwimmerImageFARCI \\
Median latency, ms & \SwimmerLatency & M4 Pro, CPU provider \\
\bottomrule
\end{tabular}
\end{table}

We do not quote the author's headline mAP@50 of \SwimmerAuthorMAP{}
(\reported{}): it is across five classes and carried by boats, and the swimmer
class — the only one VIGÍA cares about — is their second weakest at
\SwimmerAuthorAP{} (\reported{}). Our \SwimmerFOne{} and their
\SwimmerAuthorAP{} are not in conflict. Average precision over strict $\iou$
thresholds punishes imprecise boxes on tiny objects, while our
$\tau = 0.3$ operating point asks the operationally relevant question, which
is whether the person was found at all.

\begin{figure}[htbp]
\centering
\includegraphics[width=0.85\textwidth]{confidence_sweep.pdf}
\caption{Swimmer precision and recall against the confidence threshold. The
operating point is chosen where recall is still high, because the validator
can discard a false alarm but nothing downstream can recover a person never
detected — the same asymmetry that motivates $\beta = 2$ in
\cref{eq:fbeta-ratio}.}
\label{fig:sweep}
\end{figure}

\subsection{Trapped people and building access}
\label{sec:building}

\begin{description}[leftmargin=7.4em,labelwidth=7.0em,style=sameline,font=\normalfont\itshape]
  \item[Model] YOLOv8-DRN, trained and published by the DRespNeT
    authors~\citep{sirma2025drespnet}, CC BY 4.0 — attribution is legally
    required here, unlike SeaDronesSee where it is given by choice.
  \item[Dataset] DRespNeT: UAV imagery of the 2023 Türkiye earthquakes,
    650 train / 50 valid / 25 test.
  \item[Envelope] Downward and oblique UAV views of earthquake-damaged urban
    blocks. Not validated for street-level views, interior imagery or night.
    Trained on 650 images, so generalisation to other disaster types and other
    cities is unproven rather than assumed.
\end{description}

The model predicts 28 classes; \cref{eq:merge} reduces them to five at the
detector boundary, and only \texttt{civilian} reaches the validator. The model
is an instance segmenter, so its output carries 32 mask coefficients after the
class scores; these are trimmed at postprocess, because without that the
decoder reads 60 classes and takes an argmax over mask prototypes.

\begin{table}[htbp]
\centering\small
\caption{Building-access detector on the held-out test split,
\BuildingTestImages{} images, \BuildingCivilians{} ground-truth civilian
instances, confidence $0.30$. The image-level rows are computed over twelve
positive and twelve negative images; their intervals are given because the
point estimates alone would badly overstate what \BuildingTestImages{} images
can establish. \measured{} — \texttt{eval/run\_building\_access\_eval.py}.}
\label{tab:building}
\begin{tabular}{l S[table-format=1.3] l}
\toprule
Metric & {Value} & Note \\
\midrule
Civilian box precision & \BuildingBoxPrecision & \\
Civilian box recall    & \BuildingBoxRecall    & 95\,\% CI \BuildingCivilianRecallCI \\
Civilian box $F_1$     & \BuildingBoxFOne      & \\
Image-level recall     & \BuildingImageRecall  & 95\,\% CI \BuildingImageRecallCI \\
Image-level FAR        & \BuildingImageFAR     & 95\,\% CI \BuildingImageFARCI \\
Rescue team $F_1$      & \BuildingRescueFOne   & operator context, not an alarm \\
Accessible entry $F_1$ & \BuildingEntryOpenFOne & operator context \\
Blocked entry $F_1$    & \BuildingEntryBlockedFOne & operator context \\
Collapsed building $F_1$ & \BuildingCollapsedFOne & operator context \\
Median latency, ms & \BuildingLatency & M4 Pro, CPU provider \\
\bottomrule
\end{tabular}
\end{table}

The authors report 92.7\,\% mAP@50 at 27\,FPS across all 28 classes
(\reported{}). That is not comparable to \cref{tab:building}: they report over
28 classes, we merge to five and measure one.

\begin{recorded}{An architectural preference abandoned on evidence}
VIGÍA first trained its own torchvision Faster R-CNN on the same 650 images,
specifically because torchvision is BSD-3-Clause and would have kept this
detector clear of the AGPL perimeter altogether. On the identical held-out
split with the identical metric it reached civilian box recall of $0.141$ — it
missed roughly six of every seven people it should have boxed, and was by a
wide margin the weakest component in the system. The authors' own weights
reach \BuildingBoxRecall. The test/valid gap also inverted: our model scored
$0.462$ on valid against $0.141$ on test, a collapse showing it had not
generalised, where the published weights score $0.815$ against
\BuildingBoxFOne. Keeping a six-times-worse detector to preserve an
architectural preference would have been a choice against the people this
detector exists to find. The AGPL exposure is contained the same way as fire:
ONNX conversion in an isolated build environment, no weights redistributed,
enforced by \texttt{tools/check\_licence.py}.
\end{recorded}

\begin{figure}[htbp]
\centering
\includegraphics[width=0.85\textwidth]{building_access_models.pdf}
\caption{The two models on the identical held-out split with the identical
metric. Image-level recall is \BuildingImageRecall{} for both, which is
exactly why image-level metrics alone are not enough (\cref{sec:twolevels}):
they hide a six-fold difference in whether the person was actually located.}
\label{fig:building-models}
\end{figure}

\begin{figure}[htbp]
\centering
\includegraphics[width=0.85\textwidth]{building_access_per_class.pdf}
\caption{Per-class $F_1$ on the held-out split at confidence $0.30$. Only
\texttt{civilian} reaches the validator; the rest are operator context.}
\label{fig:building-classes}
\end{figure}

\subsection{Road incidents}
\label{sec:traffic}

Included to test whether the validator transfers to a hazard with entirely
different physics. The model~\citep{enostraffic} is MIT-licensed and detects
accidents and vehicles in roadside and dashcam imagery; only the
\texttt{accident} class reaches the validator, since passing \texttt{vehicle}
through would make every car on the road an alert.

Its status is \emph{not measured}. It is integrated and running behind the
validator, and it supplied the second hazard in \cref{tab:transfer}, but that
experiment measures the \emph{validator}, not this detector: no precision or
recall for the accident class against held-out labels has been run. The
available footage is a dashcam compilation that is predominantly accident
footage and therefore cannot support a false-positive rate. The authors report
mAP@50 of $0.826$ and accident-class recall of $0.855$ (\reported{}); we do
not adopt those as ours.

\begin{recorded}{A checkpoint choice recorded rather than silently made}
The author's own inference script uses the \texttt{epoch14} checkpoint. On
their four published test images, \texttt{epoch61} returns exactly one
accident per image at confidence $0.58$--$0.70$, whereas \texttt{epoch14}
returns three detections on one image. We use \texttt{epoch61} and record the
reason here rather than leaving a future reader to wonder why the two
disagree.
\end{recorded}
